The original block A contains 60 services and 300 rungs. Repeated mathematical checkpoints are byte-identical. The pilot-curvature bisection generator realizes nonuniform axes in 0 of 252 date models; its name alone does not establish effective adaptation. At common successful declared target-by-repetition comparisons, witness is faster than uniform FVI in 13 of 39 and slower in 26. Against pilot-curvature bisection FVI, the corresponding counts are 30 faster and 9 slower out of 39. These counts summarize selected targets; Table~\ref{tab:frontier49} and the exact deposited partition, rather than a binary attainment headline, are the primary work comparison.

The new first-crossing direct-cost intervals identify lower witness cost in 0 cases, higher cost in 93, and leave 3 unresolved. For the separately evaluated favorable R48 work pair, the corresponding counts are 0, 0 and 8. A cheaper first crossing and a sharper bound therefore receive an actual economic comparison rather than an inferred ranking. The detailed intervals and initial actions are retained, including every adverse outcome.

Across the 24 declared controller--price decisions, the upper-net-cost selector is certified strictly cheaper than the all-state-bound minimizer in 24. Its cost-regret upper bounds relative to the best evaluated bit range from 0.00238686 to 0.00332242 in the normalized economic units. Where the comparison does not identify a strict improvement it remains unresolved; midpoint choices are descriptive, not true optima. This is a finite expected-net-cost conclusion under the uniform initial law, not a calibrated welfare result or a universal preference for coarse observation.

The direct and sensor study contains 200 sampled contrasts, totaling 52,428,800 interval paths, with a separately recorded joint evaluation time of 923.64 seconds. Identity rows use no simulated paths. This evaluation work is not allocated to one favored generator or hidden inside another method's construction clock.
